
D'après (1.11.3), on
peut écrire
\begin{equation}\tag{1.11.4}
  \sum_{j=0}^r x_j\left(\sum_{i=0}^{r+1}(-1)^i\lambda^i(E)\sigma_{p(j)-i}\right)=0,
\end{equation}
avec \(p(j)\geq1\). On prend \(p(j)=n_0+r+j\), ce qui est bien \(\geq1\) si \(n_0>-r\). On obtient alors :
\[
\begin{aligned}
\sum_{j=0}^r x_j\left(\sum_{i=0}^{r+1}(-1)^i\lambda^i(E)\sigma_{n_0+r+j-i}\right)
&=\sum_{i=0}^{r+1}(-1)^i\lambda^i(E)
   \left(\sum_{j=0}^r x_j\sigma_{n_0+r+j-i}\right).
\end{aligned}
\]

Pour \(i\neq r+1\), \(n_0+r-i\geq n_0\), donc d'après (1.11.2) :
\[
  \sum_{j=0}^r x_j\sigma_{n_0+r+j-i}=0.
\]

Par conséquent, la relation (1.11.4) se réduit à :
\[
  \lambda^{r+1}(E).\left(\sum_{j=0}^r x_j\sigma_{n_0-1+j}\right)=0.
\]

Or \(\underline{E}\) est de rang \(r+1\), donc \(\lambda^{r+1}(E)\) est la classe d'un faisceau inversible, et est donc inversible dans \(K^{\bullet}(S)\). Par suite, la relation (1.11.4) se réduit à la relation (1.11.2) pour \(n_0-1\).

Donc la relation (1.11.2) est vraie pour \(n\geq-r\); pour \(n=-r\), elle donne :
\[
  x_r\sigma_0=0.
\]
Comme \(\sigma_0=1\), on obtient \(x_r=0\). On montre ainsi par récurrence descendante sur \(j\) que les \(x_j\) sont tous nuls, en utilisant les relations (1.11.2) pour \(n\) variant de \(-r\) à \(0\). Ceci montre que \(1,\ldots,\xi^r\) sont indépendants sur \(K^{\bullet}(S)\), et achève la démonstration du théorème 1.1.

\textbf{Remarque 1.12.} Pour tout \(\underline{O}_S\)-Module localement libre de type fini \(\underline{E}\), soit \(\sigma_i(E)\) la classe dans \(K^{\bullet}(S)\) de la \(i\)-ième puissance symétrique de \(\underline{E}\); posons \(\sigma_t(E)=1+\sum_{i\geq1}\sigma_i(E)t^i\)~; les relations (1.11.3) peuvent encore s'écrire :
